% GENERATED FILE. Edit canonical lessons, then run npm run generate:books.
% canonical-source-hash: fnv1a64:5c87fabc8fc550a0
% canonical-lessons: HI-C410-pasport, HI-C410-viza, HI-C410-suraksa, HI-C410-kastam, HI-C410-sulk

\chapter{Passports and Customs}
\label{ch:hi-passports-and-customs}
\hlchaptermodality{410}

\begin{chapteropening}
\textbf{By the end of this chapter:} \emph{I can show my passport and visa, go through security and customs, and pay a duty.}

Complete the last lesson of chapter 410: I can show my passport and visa, go through security and customs, and pay a duty.
\end{chapteropening}

\section[\texorpdfstring{\hi{पासपोर्ट} (pāsporṭ)}{पासपोर्ट (pāsporṭ)}]{\hi{पासपोर्ट} (pāsporṭ) — a passport}

\label{lesson:HI-C410-pasport}



\begin{quote}
Before the new one: say the Hindi for a flight, then the Hindi for arrival.
\end{quote}

\subsection*{You'll want to know: \hi{पासपोर्ट}}
\textbf{\hi{पासपोर्ट}} — \emph{pāsporṭ} — \textquotedblleft{}a passport\textquotedblright{}.

\begin{quote}
\hi{अपना} \hi{पासपोर्ट} \hi{हमेशा} \hi{पास} \hi{रखिए।}
\end{quote}

\begin{quote}
\emph{apnā pāsporṭ hameśā pās rakhiye.}
\end{quote}

\begin{quote}
Always keep your passport on you.
\end{quote}

Masculine.

\begin{cousinweb}
\textbf{\hi{पासपोर्ट}} (\emph{pāsporṭ}) is borrowed from English \emph{passport}.
\end{cousinweb}

\begin{grammarlens}[title={What you've built}]
One more word for getting around.
\end{grammarlens}

\subsection*{Your turn}
\begin{itemize}
\raggedright
  \item \emph{Say it:} \emph{pāsporṭ}
  \item \emph{Say it:} \emph{pāsporṭ}, once more
  \item \emph{Say it:} \emph{apnā pāsporṭ hameśā pās rakhiye.} — Always keep your passport on you.
  \item \emph{Recall:} say the Hindi for a flight, then the Hindi for arrival, then say \emph{pāsporṭ} again
\end{itemize}

\begin{tcolorbox}[breakable,colback=teal!4,colframe=teal!35!black,title={Before you move on}]
What does \hi{पासपोर्ट} mean? (\textquotedblleft{}A passport\textquotedblright{}.) Say it once more.
\end{tcolorbox}
\section[\texorpdfstring{\hi{वीज़ा} (vīzā)}{वीज़ा (vīzā)}]{\hi{वीज़ा} (vīzā) — a visa}

\label{lesson:HI-C410-viza}



\begin{quote}
Before the new one: say the Hindi for arrival, then the Hindi for a passport.
\end{quote}

\subsection*{You'll want to know: \hi{वीज़ा}}
\textbf{\hi{वीज़ा}} — \emph{vīzā} — \textquotedblleft{}a visa\textquotedblright{}.

\begin{quote}
\hi{मुझे} \hi{अभी} \hi{वीज़ा} \hi{नहीं} \hi{मिला।}
\end{quote}

\begin{quote}
\emph{mujhe abhī vīzā nahīṃ milā.}
\end{quote}

\begin{quote}
I have not got my visa yet.
\end{quote}

Masculine.

\begin{cousinweb}
\textbf{\hi{वीज़ा}} (\emph{vīzā}) is borrowed from English \emph{visa}; the English z keeps its \hi{ज़}.
\end{cousinweb}

\begin{grammarlens}[title={What you've built}]
One more word for getting around.
\end{grammarlens}

\subsection*{Your turn}
\begin{itemize}
\raggedright
  \item \emph{Say it:} \emph{vīzā}
  \item \emph{Say it:} \emph{vīzā}, once more
  \item \emph{Say it:} \emph{mujhe abhī vīzā nahīṃ milā.} — I have not got my visa yet.
  \item \emph{Recall:} say the Hindi for arrival, then the Hindi for a passport, then say \emph{vīzā} again
\end{itemize}

\begin{tcolorbox}[breakable,colback=teal!4,colframe=teal!35!black,title={Before you move on}]
What does \hi{वीज़ा} mean? (\textquotedblleft{}A visa\textquotedblright{}.) Say it once more.
\end{tcolorbox}
\section[\texorpdfstring{\hi{सुरक्षा} (surakṣā)}{सुरक्षा (surakṣā)}]{\hi{सुरक्षा} (surakṣā) — security; safety}

\label{lesson:HI-C410-suraksa}



\begin{quote}
Before the new one: say the Hindi for a passport, then the Hindi for a visa.
\end{quote}

\subsection*{You'll want to know: \hi{सुरक्षा}}
\textbf{\hi{सुरक्षा}} — \emph{surakṣā} — \textquotedblleft{}security; safety\textquotedblright{}.

\begin{quote}
\hi{एयरपोर्ट} \hi{पर} \hi{सुरक्षा} \hi{बहुत} \hi{सख़्त} \hi{थी।}
\end{quote}

\begin{quote}
\emph{eyarporṭ par surakṣā bahut sakht thī.}
\end{quote}

\begin{quote}
Security at the airport was very tight.
\end{quote}

Feminine.

\begin{cousinweb}
\textbf{\hi{सुरक्षा}} (\emph{surakṣā}) is a learned word made of Sanskrit \emph{su-}, \textquotedblleft{}well\textquotedblright{}, and \emph{rakṣā}, \textquotedblleft{}protection, guarding\textquotedblright{} (Monier-Williams). \hi{सुरक्षित}, safe, is its partner.
\end{cousinweb}

\begin{grammarlens}[title={What you've built}]
One more word for getting around.
\end{grammarlens}

\subsection*{Your turn}
\begin{itemize}
\raggedright
  \item \emph{Say it:} \emph{surakṣā}
  \item \emph{Say it:} \emph{surakṣā}, once more
  \item \emph{Say it:} \emph{eyarporṭ par surakṣā bahut sakht thī.} — Security at the airport was very tight.
  \item \emph{Recall:} say the Hindi for a passport, then the Hindi for a visa, then say \emph{surakṣā} again
\end{itemize}

\begin{tcolorbox}[breakable,colback=teal!4,colframe=teal!35!black,title={Before you move on}]
What does \hi{सुरक्षा} mean? (\textquotedblleft{}Security; safety\textquotedblright{}.) Say it once more.
\end{tcolorbox}
\section[\texorpdfstring{\hi{कस्टम} (kasṭam)}{कस्टम (kasṭam)}]{\hi{कस्टम} (kasṭam) — customs}

\label{lesson:HI-C410-kastam}



\begin{quote}
Before the new one: say the Hindi for a visa, then the Hindi for security.
\end{quote}

\subsection*{You'll want to know: \hi{कस्टम}}
\textbf{\hi{कस्टम}} — \emph{kasṭam} — \textquotedblleft{}customs\textquotedblright{}.

\begin{quote}
\hi{कस्टम} \hi{पर} \hi{एक} \hi{आदमी} \hi{ने} \hi{मेरी} \hi{अटैची} \hi{खोली।}
\end{quote}

\begin{quote}
\emph{kasṭam par ek ādmī ne merī aṭaicī kholī.}
\end{quote}

\begin{quote}
At customs a man opened my suitcase.
\end{quote}

Masculine.

\begin{cousinweb}
\textbf{\hi{कस्टम}} (\emph{kasṭam}) is borrowed from English \emph{customs}, the check on goods at an airport or a border; Hindi has dropped the final s.
\end{cousinweb}

\begin{grammarlens}[title={What you've built}]
One more word for getting around.
\end{grammarlens}

\subsection*{Your turn}
\begin{itemize}
\raggedright
  \item \emph{Say it:} \emph{kasṭam}
  \item \emph{Say it:} \emph{kasṭam}, once more
  \item \emph{Say it:} \emph{kasṭam par ek ādmī ne merī aṭaicī kholī.} — At customs a man opened my suitcase.
  \item \emph{Recall:} say the Hindi for a visa, then the Hindi for security, then say \emph{kasṭam} again
\end{itemize}

\begin{tcolorbox}[breakable,colback=teal!4,colframe=teal!35!black,title={Before you move on}]
What does \hi{कस्टम} mean? (\textquotedblleft{}Customs\textquotedblright{}.) Say it once more.
\end{tcolorbox}
\section[\texorpdfstring{\hi{शुल्क} (śulk)}{शुल्क (śulk)}]{\hi{शुल्क} (śulk) — a fee; a duty (formal)}

\label{lesson:HI-C410-sulk}



\begin{quote}
Before the new one: say the Hindi for security, then the Hindi for customs.
\end{quote}

\subsection*{You'll want to know: \hi{शुल्क}}
\textbf{\hi{शुल्क}} — \emph{śulk} — \textquotedblleft{}a fee; a duty (formal)\textquotedblright{}.

\begin{quote}
\hi{नई} \hi{घड़ी} \hi{के} \hi{लिए} \hi{हमें} \hi{कस्टम} \hi{पर} \hi{शुल्क} \hi{देना} \hi{पड़ा।}
\end{quote}

\begin{quote}
\emph{naī ghaṛī ke liye hamẽ kasṭam par śulk denā paṛā.}
\end{quote}

\begin{quote}
We had to pay duty at customs on the new watch.
\end{quote}

Masculine.

\begin{cousinweb}
\textbf{\hi{शुल्क}} (\emph{śulk}) is a tatsama, Sanskrit \emph{śulka}, \textquotedblleft{}a price\textquotedblright{} and also a toll or duty taken at roads and ferries (Monier-Williams); offices, customs and the ticket window at a monument all use it for a fee.
\end{cousinweb}

\begin{grammarlens}[title={What you've built}]
One more word for getting around.
\end{grammarlens}

\subsection*{Your turn}
\begin{itemize}
\raggedright
  \item \emph{Say it:} \emph{śulk}
  \item \emph{Say it:} \emph{śulk}, once more
  \item \emph{Say it:} \emph{naī ghaṛī ke liye hamẽ kasṭam par śulk denā paṛā.} — We had to pay duty at customs on the new watch.
  \item \emph{Recall:} say the Hindi for security, then the Hindi for customs, then say \emph{śulk} again
\end{itemize}

\begin{tcolorbox}[breakable,colback=teal!4,colframe=teal!35!black,title={Before you move on}]
What does \hi{शुल्क} mean? (\textquotedblleft{}A fee; a duty (formal)\textquotedblright{}.) Say it once more.
\end{tcolorbox}
